% =====================================================================
\chapter{Funções e gráficos}

O Cálculo é a matemática que descreve como as funções variam. Neste capítulo revisamos
todas as funções necessárias ao estudo do Cálculo: polinomiais, racionais, trigonométricas,
exponenciais, logarítmicas e hiperbólicas. Vemos como avaliá-las, as propriedades dos seus
gráficos e as técnicas algébricas para resolver equações que as envolvem. É a base de tudo o
que vem depois; convém dominá-la antes da introdução formal do Cálculo, no próximo capítulo.

% ---------------------------------------------------------------------
\section{Revisão de funções}

\subsection{Funções}

Dados dois conjuntos $A$ e $B$, um conjunto de pares ordenados $(x,y)$ com $x\in A$ e
$y\in B$ é uma \emph{relação} de $A$ em $B$. Uma função é um tipo especial de relação, em
que cada elemento do primeiro conjunto está associado a \emph{exatamente um} elemento do
segundo. O elemento do primeiro conjunto é a \emph{entrada}; o do segundo, a \emph{saída}.
Quando conhecemos a entrada, a saída fica determinada; por isso dizemos que a saída é
função da entrada. A área de um quadrado, por exemplo, é função do comprimento do lado.

\begin{definicao}
Uma \textbf{função}\index{funcao@função} $f$ consiste em um conjunto de entradas, um conjunto de saídas e uma
regra que associa cada entrada a exatamente uma saída. O conjunto de entradas é o
\textbf{domínio}\index{dominio@domínio} da função; o conjunto de saídas é a \textbf{imagem}\index{imagem}.
\end{definicao}

Para uma função $f$ com domínio $D$, costumamos usar $x$ para a entrada e $y$ para a saída,
e escrevemos $y=f(x)$ (lê-se ``$y$ é igual a $f$ de $x$''). Chamamos $x$ de \emph{variável
independente} e $y$ de \emph{variável dependente}. Uma entrada só pode ser enviada a uma
saída, mas duas entradas diferentes podem ser enviadas à mesma saída
(Figura~\thechapter.1).

\figtikz{\begin{tikzpicture}[font=\small]
\draw[rounded corners=8pt,fill=azul!10,draw=azul] (0,0) rectangle (2,3.6);
\draw[rounded corners=8pt,fill=azul!10,draw=azul] (4.2,0) rectangle (6.2,3.6);
\node at (1,-0.35) {Domínio}; \node at (5.2,-0.35) {Imagem};
\foreach \y/\t in {3/1,2.2/2,1.4/3,0.6/4} \node (a\t) at (1,\y) {\t};
\foreach \y/\t in {2.8/2,1.8/4,0.8/6} \node (b\t) at (5.2,\y) {\t};
\draw[-{Stealth},laranja,thick] (a1) -- (b6);
\draw[-{Stealth},laranja,thick] (a2) -- (b4);
\draw[-{Stealth},azul,thick] (a3) -- (b2);
\draw[-{Stealth},red!70!black,thick] (a4) -- (b2);
\draw[-{Stealth},thick] (1.6,3.9) to[bend left=20] node[above]{$f$} (4.6,3.9);
\end{tikzpicture}}{Uma função associa cada elemento do domínio a exatamente um elemento
da imagem. Entradas diferentes (aqui, 3 e 4) podem ter a mesma saída.}

O \textbf{gráfico}\index{grafico de uma funcao@gráfico de uma função} de $f$ é o conjunto dos pontos $(x,f(x))$ do plano, com $x$ no domínio.
Quando uma função é dada apenas por uma fórmula, sem domínio explícito, o domínio é o
conjunto de todos os reais $x$ para os quais $f(x)$ é um número real; é o chamado
\emph{domínio natural}. Assim, o domínio de $f(x)=x^2$ é $\R$, e o de $f(x)=\sqrt{x}$ é o
conjunto dos reais não negativos.

\textbf{Notação.} Usamos a notação de conjuntos, como $\Set{x}{1<x<5}$, e a de intervalos:
\[
(1,5)=\Set{x}{1<x<5},\quad [1,5]=\Set{x}{1\le x\le 5},\quad
[0,\infty)=\Set{x}{x\ge 0},\quad (-\infty,\infty)=\R .
\]
O símbolo $\infty$ não é um número real; indica apenas que o conjunto não tem limitação
naquela direção.

Algumas funções são \textbf{definidas por partes}\index{funcao@função!definida por partes}, com fórmulas diferentes em diferentes
partes do domínio. Por exemplo,
\[
f(x)=\begin{cases}3x+1, & x\ge 2\\ x^2, & x<2\end{cases}
\]
dá $f(5)=3(5)+1=16$ e $f(-1)=(-1)^2=1$.

\begin{exemplo}{1.1}{Determinando domínio e imagem}
Determine o domínio e a imagem de (a) $f(x)=(x-4)^2+5$; (b) $f(x)=\sqrt{3x+2}-1$;
(c) $f(x)=\dfrac{3}{x-2}$.
\solucao
(a) $f(x)$ é real para todo $x$, logo o domínio é $(-\infty,\infty)$. Como $(x-4)^2\ge0$,
temos $f(x)\ge5$. Para todo $y\ge5$, a equação $(x-4)^2+5=y$ tem solução
$x=4\pm\sqrt{y-5}$; logo a imagem é $\Set{y}{y\ge5}$.

(b) Precisamos de $3x+2\ge0$, ou seja, domínio $\Set{x}{x\ge -2/3}$. Como
$\sqrt{3x+2}\ge0$, temos $f(x)\ge-1$. Dado $y\ge-1$, de $\sqrt{3x+2}=y+1$ vem
$x=\tfrac13(y+1)^2-\tfrac23\ge-\tfrac23$, que está no domínio. A imagem é $\Set{y}{y\ge-1}$.

(c) O denominador não pode se anular: domínio $\Set{x}{x\neq2}$. De $3/(x-2)=y$ vem
$x=3/y+2$, possível sempre que $y\neq0$. A imagem é $\Set{y}{y\neq0}$.
\end{exemplo}

\subsection{Representações de uma função}

Uma função pode ser representada por uma \emph{tabela}, um \emph{gráfico} ou uma
\emph{fórmula}, e muitas vezes usamos as três juntas: marcamos no gráfico os valores de uma
tabela, ou construímos uma tabela a partir de uma fórmula.

Os valores de $x$ com $f(x)=0$ são os \textbf{zeros}\index{zeros de uma funcao@zeros de uma função} da função; eles indicam onde o gráfico
corta o eixo $x$. O \textbf{intercepto $y$}, se existir, é o ponto $(0,f(0))$. Como cada
entrada tem uma única saída, cada reta vertical $x=c$ corta o gráfico no máximo uma vez.

\begin{regra}{Regra: teste da reta vertical}\index{teste da reta vertical}
Dada uma função $f$, toda reta vertical corta o gráfico de $f$ no máximo uma vez. Se alguma
reta vertical corta um conjunto de pontos mais de uma vez, esse conjunto não é gráfico de
uma função.
\end{regra}

\figtikz{\begin{tikzpicture}
\begin{groupplot}[group style={group size=2 by 1,horizontal sep=1.2cm},painel,width=5.6cm,height=4.8cm,xtick=\empty,ytick=\empty]
\nextgroupplot[xmin=-2.4,xmax=2.6,ymin=-2.2,ymax=2.4,title={(a) é função}]
\foreach \a in {-1.4,0.3,1.6}{\edef\tmp{\noexpand\draw[laranja,thin] (axis cs:\a,-2.2) -- (axis cs:\a,2.4);}\tmp}
\addplot[domain=-2.2:2.4]{0.25*x^4-0.15*x^3-1.1*x^2+0.4*x+0.3};
\nextgroupplot[xmin=-2.2,xmax=2.2,ymin=-2.2,ymax=2.4,title={(b) não é função}]
\foreach \a in {-0.6,0.5}{\edef\tmp{\noexpand\draw[laranja,thin] (axis cs:\a,-2.2) -- (axis cs:\a,2.4);}\tmp}
\addplot[variable=t,domain=-1.75:1.75] ({t^3-2*t},{t});
\end{groupplot}\end{tikzpicture}}{(a) Gráfico de uma função: toda reta vertical o
corta no máximo uma vez. (b) Não é gráfico de função: algumas retas verticais o cortam mais
de uma vez.}

Dizemos que $f$ é \textbf{crescente}\index{funcao@função!crescente} em um intervalo $I$ se $f(x_1)\le f(x_2)$ sempre que
$x_1<x_2$ em $I$, e \textbf{estritamente crescente} se $f(x_1)<f(x_2)$. Analogamente, $f$
é \textbf{decrescente}\index{funcao@função!decrescente} se $f(x_1)\ge f(x_2)$ e \textbf{estritamente decrescente} se
$f(x_1)>f(x_2)$ sempre que $x_1<x_2$. Por exemplo, $f(x)=3x$ é crescente em $\R$, e
$f(x)=-x^3$ é estritamente decrescente em $\R$.

\subsection{Combinação e composição de funções}

Dadas $f$ e $g$, definimos a soma, a diferença, o produto e o quociente:
\[
(f\pm g)(x)=f(x)\pm g(x),\qquad (f\cdot g)(x)=f(x)\,g(x),\qquad
\Big(\frac fg\Big)(x)=\frac{f(x)}{g(x)}\ \ \text{se } g(x)\neq0 .
\]
Por exemplo, com $f(x)=2x-3$ e $g(x)=x^2-1$, temos $(f/g)(x)=\dfrac{2x-3}{x^2-1}$, com
domínio $\Set{x}{x\neq\pm1}$.

Também podemos aplicar uma função ao resultado de outra. Se o custo $C$ de aquecer um
prédio depende da temperatura $T$, e a temperatura depende da hora $t$, então o custo como
função do tempo é $C(T(t))$.

\begin{definicao}
Sejam $f$ uma função com domínio $A$ e imagem $B$, e $g$ uma função com domínio $D$ e
imagem $E$. Se $B\subseteq D$, a \textbf{função composta}\index{funcao@função!composta} $g\circ f$ é a função com
domínio $A$ dada por
\[
(g\circ f)(x)=g\big(f(x)\big).
\]
\end{definicao}

\begin{exemplo}{1.2}{Composição de funções dadas por fórmulas}
Sejam $f(x)=x^2+1$ e $g(x)=1/x$. Encontre $g\circ f$ e $f\circ g$, com domínio e imagem.
\solucao
$(g\circ f)(x)=g(x^2+1)=\dfrac{1}{x^2+1}$. Como $x^2+1\neq0$ para todo $x$, o domínio é
$\R$. Como $0<1/(x^2+1)\le1$, e para cada $y\in(0,1]$ a equação tem solução
$x=\pm\sqrt{1/y-1}$, a imagem é $(0,1]$.

$(f\circ g)(x)=f(1/x)=\dfrac{1}{x^2}+1$, com domínio $\Set{x}{x\neq0}$. Resolvendo
$(1/x)^2+1=y$, obtemos $x=\pm1/\sqrt{y-1}$, real se e somente se $y>1$. A imagem é
$\Set{y}{y>1}$.

Note que $f\circ g\neq g\circ f$: a ordem da composição importa.
\end{exemplo}

\subsection{Simetria}

Se $f(-x)=f(x)$, o ponto $(-x,y)$ está no gráfico sempre que $(x,y)$ está, e o gráfico é
simétrico em relação ao eixo $y$. Se $f(-x)=-f(x)$, o gráfico é simétrico em relação à
origem: girá-lo $180^\circ$ em torno da origem não o altera.

\begin{definicao}
Se $f(-x)=f(x)$ para todo $x$ do domínio, $f$ é uma \textbf{função par}\index{funcao@função!par} (simétrica em
relação ao eixo $y$). Se $f(-x)=-f(x)$ para todo $x$ do domínio, $f$ é uma
\textbf{função ímpar}\index{funcao@função!impar@ímpar} (simétrica em relação à origem).
\end{definicao}

\figtikz{\begin{tikzpicture}
\begin{groupplot}[group style={group size=2 by 1,horizontal sep=1.2cm},painel,width=5.6cm,height=4.8cm]
\nextgroupplot[xmin=-2.6,xmax=2.6,ymin=-5,ymax=5,xtick={-2,-1,1,2},ytick={-4,-3,2,4},title={(a) $x^4-2x^2-3$: par}]
\addplot[domain=-2.25:2.25]{x^4-2*x^2-3};
\nextgroupplot[xmin=-3,xmax=3,ymin=-5,ymax=5,xtick={-2,2},ytick={-3,3},title={(b) $x^3-4x$: ímpar}]
\addplot[domain=-2.75:2.75]{x^3-4*x};
\end{groupplot}\end{tikzpicture}}{(a) $f(x)=x^4-2x^2-3$ é par: simetria em relação
ao eixo $y$. (b) $f(x)=x^3-4x$ é ímpar: simetria em relação à origem.}

\begin{exemplo}{1.3}{Funções pares e ímpares}
Classifique: (a) $f(x)=-5x^4+7x^2-2$; (b) $f(x)=2x^5-4x+5$; (c) $f(x)=\dfrac{3x}{x^2+1}$.
\solucao
Calculamos $f(-x)$ e comparamos com $f(x)$ e $-f(x)$.
(a) $f(-x)=-5x^4+7x^2-2=f(x)$: par.
(b) $f(-x)=-2x^5+4x+5$, que não é $f(x)$ nem $-f(x)=-2x^5+4x-5$: nem par nem ímpar.
(c) $f(-x)=\dfrac{-3x}{x^2+1}=-f(x)$: ímpar.
\end{exemplo}

A \textbf{função valor absoluto}\index{valor absoluto} é
\[
|x|=\begin{cases}-x, & x<0\\ x, & x\ge0.\end{cases}
\]
Sua imagem é $\Set{y}{y\ge0}$, e ela é par. Por exemplo, $f(x)=2|x-3|+4$ tem domínio $\R$
e, como $|x-3|\ge0$, imagem $\Set{y}{y\ge4}$.

% ---------------------------------------------------------------------
\section{Classes básicas de funções}

\subsection{Funções afins e inclinação}

Uma \textbf{função afim}\index{funcao@função!afim} (ou linear) tem a forma $f(x)=mx+b$, e seu gráfico é uma reta.
Para quaisquer dois pontos $(x_1,y_1)$ e $(x_2,y_2)$ dessa reta, a razão entre a variação
de $y$ e a de $x$ é sempre a mesma.

\begin{definicao}
A \textbf{inclinação}\index{inclinacao@inclinação} da reta que passa por $(x_1,y_1)$ e $(x_2,y_2)$ é
\[
m=\frac{y_2-y_1}{x_2-x_1}=\frac{\Delta y}{\Delta x}.
\]
A reta que passa por $(x_1,y_1)$ com inclinação $m$ tem \textbf{equação ponto-inclinação}\index{equacao da reta@equação da reta}
$y-y_1=m(x-x_1)$. A reta com inclinação $m$ e intercepto $(0,b)$ tem equação
$y=mx+b$ (forma \textbf{reduzida}). A \textbf{forma geral} $ax+by=c$, com $a$ e $b$ não
ambos nulos, inclui também as retas verticais $x=k$.
\end{definicao}

\begin{exemplo}{1.4}{Inclinação e equação de uma reta}
Encontre a reta que passa por $(11,-4)$ e $(-4,5)$.
\solucao
$m=\dfrac{5-(-4)}{-4-11}=-\dfrac{9}{15}=-\dfrac35$. Pelo ponto $(11,-4)$:
$y+4=-\tfrac35(x-11)$, ou seja, $y=-\tfrac35x+\tfrac{13}{5}$.
\end{exemplo}

\subsection{Polinômios}

Um \textbf{polinômio}\index{polinomio@polinômio} é uma função da forma
$f(x)=a_nx^n+a_{n-1}x^{n-1}+\dots+a_1x+a_0$, com $n$ inteiro não negativo e $a_n\neq0$; $n$
é o \emph{grau}. Polinômios de grau 0, 1, 2 e 3 são chamados constantes, afins,
quadráticos e cúbicos. As \textbf{funções potência}\index{funcao@função!potencia@potência} $f(x)=ax^b$ incluem $x^n$, que é par
para $n$ par e ímpar para $n$ ímpar.

\figtikz{\begin{tikzpicture}
\begin{groupplot}[group style={group size=2 by 1,horizontal sep=1.2cm},painel,width=5.6cm,height=4.8cm,xmin=-1.7,xmax=1.7,xtick={-1,1}]
\nextgroupplot[ymin=-0.4,ymax=3.2,ytick={1,2},title={(a) $n$ par}]
\addplot[laranja,domain=-1.7:1.7]{x^2} node[pos=0.97,left,font=\scriptsize]{$x^2$};
\addplot[domain=-1.33:1.33]{x^4} node[pos=0.98,right,font=\scriptsize]{$x^4$};
\nextgroupplot[ymin=-3.2,ymax=3.2,ytick={-1,1},title={(b) $n$ ímpar}]
\addplot[laranja,domain=-1.47:1.47]{x^3} node[pos=0.98,right,font=\scriptsize]{$x^3$};
\addplot[domain=-1.26:1.26]{x^5} node[pos=0.98,left,font=\scriptsize]{$x^5$};
\end{groupplot}\end{tikzpicture}}{(a) Para $n$ par, $ax^n$ é par. (b) Para $n$
ímpar, $ax^n$ é ímpar.}

\textbf{Comportamento no infinito.} Para $|x|$ grande, o termo de maior grau domina. Se
$a_n>0$ e $n$ é par, $f(x)\to\infty$ quando $x\to\pm\infty$; se $n$ é ímpar,
$f(x)\to\infty$ quando $x\to\infty$ e $f(x)\to-\infty$ quando $x\to-\infty$. Se $a_n<0$,
os sinais se invertem.

\textbf{Zeros.} Os zeros de polinômios de grau 1 e 2 se encontram diretamente; para grau
maior, fatoramos quando possível.

\begin{regra}{Regra: fórmula quadrática}\index{formula quadratica@fórmula quadrática}
As soluções de $ax^2+bx+c=0$, com $a\neq0$, são
\[
x=\frac{-b\pm\sqrt{b^2-4ac}}{2a}.
\]
Se o discriminante $b^2-4ac$ é positivo, há duas soluções reais; se é nulo, uma; se é
negativo, nenhuma.
\end{regra}

\subsection{Funções algébricas e transcendentes}

Uma \textbf{função racional}\index{funcao@função!racional} é um quociente de polinômios, $f(x)=p(x)/q(x)$, definida onde
$q(x)\neq0$. As \textbf{funções algébricas}\index{funcao@função!algebrica@algébrica} são as obtidas de polinômios por soma,
subtração, multiplicação, divisão e extração de raízes; por exemplo,
$f(x)=\sqrt{x^2-4}$ e $g(x)=\sqrt[3]{x}/(1-x^2)$. Para $f(x)=\sqrt[n]{x}$, o domínio é
$[0,\infty)$ se $n$ é par e $\R$ se $n$ é ímpar. As funções que não são algébricas são as
\textbf{transcendentes}\index{funcao@função!transcendente}: trigonométricas, exponenciais, logarítmicas e hiperbólicas.

\subsection{Funções definidas por partes}

\begin{exemplo}{1.5}{Gráfico de uma função definida por partes}
Esboce $f(x)=\begin{cases}x+3, & x<1\\ (x-2)^2, & x\ge1.\end{cases}$
\solucao
Desenhamos a reta $y=x+3$ em $(-\infty,1)$ e a parábola $y=(x-2)^2$ em $[1,\infty)$. Como
$f(1)=(1-2)^2=1$, marcamos uma bola cheia em $(1,1)$. A reta se aproxima de $(1,4)$ sem
atingi-lo: marcamos ali uma bola vazia.
\figtikz{\begin{tikzpicture}
\begin{axis}[painel,width=6.2cm,height=5cm,xmin=-4.5,xmax=4.8,ymin=-1.8,ymax=6,xtick={-3,1,2,4},ytick={1,4}]
\addplot[domain=-4.5:1]{x+3} node[pos=0.25,above left,font=\scriptsize]{$x+3$};
\addplot[laranja,domain=1:4.4]{(x-2)^2} node[pos=0.9,right,font=\scriptsize]{$(x-2)^2$};
\draw[azul,fill=white] (axis cs:1,4) circle (2pt); \fill[laranja] (axis cs:1,1) circle (2pt);
\end{axis}\end{tikzpicture}}{Função afim para $x<1$ e quadrática para $x\ge1$.}
\end{exemplo}

\subsection{Transformações de funções}\index{transformacoes de graficos@transformações de gráficos}

A partir do gráfico de $y=f(x)$ obtemos outros sem novos cálculos. Para $c>0$:

\begin{center}
\begin{tabular}{@{}ll@{}}\toprule
\textbf{Função} & \textbf{Efeito sobre o gráfico de $f$}\\\midrule
$f(x)+c$ / $f(x)-c$ & translação vertical de $c$ para cima / para baixo\\
$f(x+c)$ / $f(x-c)$ & translação horizontal de $c$ para a esquerda / para a direita\\
$c\,f(x)$ & alongamento vertical se $c>1$; compressão vertical se $0<c<1$\\
$f(cx)$ & compressão horizontal se $c>1$; alongamento horizontal se $0<c<1$\\
$-f(x)$ & reflexão em relação ao eixo $x$\\
$f(-x)$ & reflexão em relação ao eixo $y$\\\bottomrule
\end{tabular}
\end{center}

\figtikz{\begin{tikzpicture}
\begin{groupplot}[group style={group size=2 by 1,horizontal sep=1.4cm},painel,width=6.2cm,height=5cm]
\nextgroupplot[xmin=-4,xmax=4,ymin=-5.5,ymax=9,title={Translações verticais}]
\addplot[cinza,dashed,domain=-4:4]{x^2};
\addplot[domain=-4:4]{x^2+4} node[pos=0.85,right,font=\scriptsize]{$x^2+4$};
\addplot[laranja,domain=-4:4]{x^2-4} node[pos=0.93,right,font=\scriptsize]{$x^2-4$};
\nextgroupplot[xmin=-7,xmax=7,ymin=-1,ymax=6,title={Translações horizontais}]
\addplot[cinza,dashed,domain=-7:7]{abs(x)};
\addplot[domain=-7:7]{abs(x+3)} node[pos=0.04,above right,font=\scriptsize]{$|x+3|$};
\addplot[laranja,domain=-7:7]{abs(x-3)} node[pos=0.96,above left,font=\scriptsize]{$|x-3|$};
\end{groupplot}\end{tikzpicture}}{Translações dos gráficos de $x^2$ e $|x|$ (tracejados).
À esquerda, verticais: $x^2\pm4$. À direita, horizontais: $|x+3|$ desloca para a esquerda
e $|x-3|$, para a direita.}

\begin{exemplo}{1.6}{Transformando uma função}
Esboce (a) $f(x)=-|x+2|-3$ e (b) $f(x)=3\sqrt{-x}+1$.
\solucao
(a) Partindo de $y=|x|$: translação de 2 unidades para a esquerda, reflexão em relação ao
eixo $x$ e translação de 3 unidades para baixo.
(b) Partindo de $y=\sqrt{x}$: reflexão em relação ao eixo $y$, alongamento vertical por um
fator 3 e translação de 1 unidade para cima.
\figtikz{\begin{tikzpicture}
\begin{groupplot}[group style={group size=2 by 1,horizontal sep=1.2cm},painel,width=5.8cm,height=5cm]
\nextgroupplot[xmin=-6.5,xmax=4,ymin=-7.5,ymax=4,xtick={-2,2},ytick={-3,3},title={(a) $-|x+2|-3$}]
\addplot[cinza,dashed,domain=-4:4]{abs(x)} node[pos=0.95,left,font=\scriptsize]{$|x|$};
\addplot[domain=-6.5:2.5]{-abs(x+2)-3};
\fill[azul] (axis cs:-2,-3) circle (1.6pt);
\nextgroupplot[xmin=-6.5,xmax=4.5,ymin=-0.8,ymax=9,xtick={-4,-1,1,4},ytick={1,4,7},title={(b) $3\sqrt{-x}+1$}]
\addplot[cinza,dashed,domain=0:4.5,samples=100]{sqrt(x)} node[pos=0.9,above,font=\scriptsize]{$\sqrt x$};
\addplot[domain=-6.5:0,samples=100]{3*sqrt(-x)+1};
\fill[azul] (axis cs:0,1) circle (1.6pt); \fill[azul] (axis cs:-4,7) circle (1.6pt);
\end{groupplot}\end{tikzpicture}}{As funções (a)
e (b) como sequências de três transformações de $|x|$ e de $\sqrt x$.}
\end{exemplo}

% ---------------------------------------------------------------------
\section{Funções trigonométricas}\index{funcoes trigonometricas@funções trigonométricas}

\subsection{Radianos}

No Cálculo, os ângulos são medidos em \textbf{radianos}\index{radiano}. A medida em radianos de um ângulo
$\theta$ é o comprimento $s$ do arco que ele determina no círculo unitário. Como a
circunferência unitária mede $2\pi$, temos $2\pi$ rad $=360^\circ$, ou seja,
\[
1\ \text{rad}=\frac{180^\circ}{\pi}\approx57{,}3^\circ,\qquad 1^\circ=\frac{\pi}{180}\ \text{rad}.
\]
Assim, $225^\circ=225\cdot\pi/180=5\pi/4$ rad, e $5\pi/3$ rad $=300^\circ$.

\subsection{As seis funções trigonométricas}

\begin{definicao}
Seja $P=(x,y)$ um ponto do círculo unitário centrado na origem $O$, e $\theta$ o ângulo
com lado inicial no semieixo $x$ positivo e lado final $OP$. Então
\[
\sen\theta=y,\quad \cos\theta=x,\quad \tg\theta=\frac yx,\quad
\cossec\theta=\frac1y,\quad \sec\theta=\frac1x,\quad \cotg\theta=\frac xy .
\]
Se $x=0$, $\sec\theta$ e $\tg\theta$ não estão definidas; se $y=0$, $\cotg\theta$ e
$\cossec\theta$ não estão definidas.
\end{definicao}

\figtikz{\begin{tikzpicture}[scale=1.9,font=\small]
\draw[->] (-1.3,0) -- (1.4,0) node[right]{$x$}; \draw[->] (0,-1.3) -- (0,1.4) node[above]{$y$};
\draw[thick] (0,0) circle (1);
\draw[thick,azul] (0,0) -- node[above left,pos=0.55]{$1$} (60:1); \fill[azul] (60:1) circle (1.3pt) node[above right]{$P=(\cos\theta,\sen\theta)$};
\draw[laranja,dashed] (60:1) -- (0.5,0) node[below]{$\cos\theta$};
\draw[->] (0.3,0) arc (0:60:0.3); \node at (25:0.45) {$\theta$};
\draw[laranja,thin] (0.5,0.45) -- (1.08,0.62) node[right,font=\footnotesize]{$\sen\theta$};
\end{tikzpicture}}{O ângulo $\theta$ na posição padrão: o ponto $P$ do círculo unitário tem coordenadas $(\cos\theta,\sen\theta)$.}

Em um círculo de raio $r$, as coordenadas de um ponto satisfazem $x=r\cos\theta$ e
$y=r\sen\theta$. Em um triângulo retângulo com ângulo agudo $\theta$, isso dá as razões
conhecidas: $\sen\theta=$ cateto oposto/hipotenusa, $\cos\theta=$ cateto
adjacente/hipotenusa e $\tg\theta=$ oposto/adjacente.

\begin{center}
\renewcommand{\arraystretch}{1.3}
\begin{tabular}{@{}c|ccccc@{}}\toprule
$\theta$ & $0$ & $\pi/6$ & $\pi/4$ & $\pi/3$ & $\pi/2$\\\midrule
$\sen\theta$ & $0$ & $1/2$ & $\sqrt2/2$ & $\sqrt3/2$ & $1$\\
$\cos\theta$ & $1$ & $\sqrt3/2$ & $\sqrt2/2$ & $1/2$ & $0$\\
$\tg\theta$ & $0$ & $\sqrt3/3$ & $1$ & $\sqrt3$ & não definida\\\bottomrule
\end{tabular}
\end{center}

\subsection{Identidades trigonométricas}

\begin{regra}{Regra: identidades trigonométricas}\index{identidades trigonometricas@identidades trigonométricas}
\textbf{Recíprocas:} $\cotg\theta=\dfrac1{\tg\theta}$,\; $\cossec\theta=\dfrac1{\sen\theta}$,\;
$\sec\theta=\dfrac1{\cos\theta}$.

\textbf{De razão:} $\tg\theta=\dfrac{\sen\theta}{\cos\theta}$,\;
$\cotg\theta=\dfrac{\cos\theta}{\sen\theta}$.

\textbf{Pitagóricas:} $\sen^2\theta+\cos^2\theta=1$,\quad $1+\tg^2\theta=\sec^2\theta$,\quad
$1+\cotg^2\theta=\cossec^2\theta$.

\textbf{Adição e subtração:}
$\sen(\alpha\pm\beta)=\sen\alpha\cos\beta\pm\cos\alpha\sen\beta$,\quad
$\cos(\alpha\pm\beta)=\cos\alpha\cos\beta\mp\sen\alpha\sen\beta$.

\textbf{Ângulo duplo:} $\sen(2\theta)=2\sen\theta\cos\theta$,\quad
$\cos(2\theta)=2\cos^2\theta-1=1-2\sen^2\theta=\cos^2\theta-\sen^2\theta$.
\end{regra}

\begin{exemplo}{1.7}{Resolvendo equações trigonométricas}
Encontre todas as soluções de (a) $1+\cos(2\theta)=\cos\theta$ e (b) $\sen(2\theta)=\tg\theta$.
\solucao
(a) Pela fórmula do ângulo duplo, a equação equivale a $2\cos^2\theta-\cos\theta=0$.
Devemos \emph{fatorar}, e não dividir por $\cos\theta$, que pode ser zero:
$\cos\theta\,(2\cos\theta-1)=0$. Logo $\cos\theta=0$ ou $\cos\theta=1/2$, e as soluções são
\[
\theta=\frac\pi2+n\pi,\qquad \theta=\pm\frac\pi3+2n\pi,\qquad n\in\mathbb Z.
\]
(b) Escrevendo $2\sen\theta\cos\theta=\sen\theta/\cos\theta$ e multiplicando por
$\cos\theta$, obtemos $\sen\theta\,(2\cos^2\theta-1)=0$. Como multiplicamos por
$\cos\theta$, que poderia ser nulo, é preciso descartar soluções estranhas no fim. Temos
$\sen\theta=0$ ou $\cos^2\theta=1/2$, e as soluções são
$\theta=n\pi$ e $\theta=\dfrac\pi4+\dfrac{n\pi}2$, $n\in\mathbb Z$.
\end{exemplo}

\subsection{Gráficos e períodos}

Uma função $f$ é \textbf{periódica}\index{funcao@função!periodica@periódica} se existe $P>0$ tal que $f(x+P)=f(x)$ para todo $x$; o
menor desses $P$ é o \textbf{período}. Seno, cosseno, secante e cossecante têm período
$2\pi$; tangente e cotangente têm período $\pi$.

\figtikz{\begin{tikzpicture}
\begin{groupplot}[group style={group size=2 by 3,horizontal sep=1.3cm,vertical sep=1.3cm},painel,width=6cm,height=3.6cm]
\nextgroupplot[title={$y=\sen x$},xmin=-6.6,xmax=6.6,ymin=-1.6,ymax=1.6,xtick={-6.2832,-3.1416,3.1416,6.2832},xticklabels={$-2\pi$,$-\pi$,$\pi$,$2\pi$}]
\addplot[domain=-6.6:6.6]{sin(deg(x))};\nextgroupplot[title={$y=\cos x$},xmin=-6.6,xmax=6.6,ymin=-1.6,ymax=1.6,xtick={-6.2832,-3.1416,3.1416,6.2832},xticklabels={$-2\pi$,$-\pi$,$\pi$,$2\pi$}]
\addplot[domain=-6.6:6.6]{cos(deg(x))};\nextgroupplot[title={$y=\tg x$},xmin=-6.6,xmax=6.6,ymin=-4,ymax=4,xtick={-6.2832,-3.1416,3.1416,6.2832},xticklabels={$-2\pi$,$-\pi$,$\pi$,$2\pi$}]
\addplot[domain=-6.560:-4.752,restrict y to domain=-6:6]{tan(deg(x))};
\addplot[domain=-4.672:-1.611,restrict y to domain=-6:6]{tan(deg(x))};
\addplot[domain=-1.531:1.531,restrict y to domain=-6:6]{tan(deg(x))};
\addplot[domain=1.611:4.672,restrict y to domain=-6:6]{tan(deg(x))};
\addplot[domain=4.752:6.560,restrict y to domain=-6:6]{tan(deg(x))};
\draw[assint] (axis cs:-4.7124,-4) -- (axis cs:-4.7124,4);
\draw[assint] (axis cs:-1.5708,-4) -- (axis cs:-1.5708,4);
\draw[assint] (axis cs:1.5708,-4) -- (axis cs:1.5708,4);
\draw[assint] (axis cs:4.7124,-4) -- (axis cs:4.7124,4);
\nextgroupplot[title={$y=\cotg x$},xmin=-6.6,xmax=6.6,ymin=-4,ymax=4,xtick={-6.2832,-3.1416,3.1416,6.2832},xticklabels={$-2\pi$,$-\pi$,$\pi$,$2\pi$}]
\addplot[domain=-6.560:-6.323,restrict y to domain=-6:6]{cos(deg(x))/sin(deg(x))};
\addplot[domain=-6.243:-3.182,restrict y to domain=-6:6]{cos(deg(x))/sin(deg(x))};
\addplot[domain=-3.102:-0.040,restrict y to domain=-6:6]{cos(deg(x))/sin(deg(x))};
\addplot[domain=0.040:3.102,restrict y to domain=-6:6]{cos(deg(x))/sin(deg(x))};
\addplot[domain=3.182:6.243,restrict y to domain=-6:6]{cos(deg(x))/sin(deg(x))};
\addplot[domain=6.323:6.560,restrict y to domain=-6:6]{cos(deg(x))/sin(deg(x))};
\draw[assint] (axis cs:-6.2832,-4) -- (axis cs:-6.2832,4);
\draw[assint] (axis cs:-3.1416,-4) -- (axis cs:-3.1416,4);
\draw[assint] (axis cs:0.0000,-4) -- (axis cs:0.0000,4);
\draw[assint] (axis cs:3.1416,-4) -- (axis cs:3.1416,4);
\draw[assint] (axis cs:6.2832,-4) -- (axis cs:6.2832,4);
\nextgroupplot[title={$y=\sec x$},xmin=-6.6,xmax=6.6,ymin=-4,ymax=4,xtick={-6.2832,-3.1416,3.1416,6.2832},xticklabels={$-2\pi$,$-\pi$,$\pi$,$2\pi$}]
\addplot[domain=-6.560:-4.752,restrict y to domain=-6:6]{1/cos(deg(x))};
\addplot[domain=-4.672:-1.611,restrict y to domain=-6:6]{1/cos(deg(x))};
\addplot[domain=-1.531:1.531,restrict y to domain=-6:6]{1/cos(deg(x))};
\addplot[domain=1.611:4.672,restrict y to domain=-6:6]{1/cos(deg(x))};
\addplot[domain=4.752:6.560,restrict y to domain=-6:6]{1/cos(deg(x))};
\draw[assint] (axis cs:-4.7124,-4) -- (axis cs:-4.7124,4);
\draw[assint] (axis cs:-1.5708,-4) -- (axis cs:-1.5708,4);
\draw[assint] (axis cs:1.5708,-4) -- (axis cs:1.5708,4);
\draw[assint] (axis cs:4.7124,-4) -- (axis cs:4.7124,4);
\nextgroupplot[title={$y=\cossec x$},xmin=-6.6,xmax=6.6,ymin=-4,ymax=4,xtick={-6.2832,-3.1416,3.1416,6.2832},xticklabels={$-2\pi$,$-\pi$,$\pi$,$2\pi$}]
\addplot[domain=-6.560:-6.323,restrict y to domain=-6:6]{1/sin(deg(x))};
\addplot[domain=-6.243:-3.182,restrict y to domain=-6:6]{1/sin(deg(x))};
\addplot[domain=-3.102:-0.040,restrict y to domain=-6:6]{1/sin(deg(x))};
\addplot[domain=0.040:3.102,restrict y to domain=-6:6]{1/sin(deg(x))};
\addplot[domain=3.182:6.243,restrict y to domain=-6:6]{1/sin(deg(x))};
\addplot[domain=6.323:6.560,restrict y to domain=-6:6]{1/sin(deg(x))};
\draw[assint] (axis cs:-6.2832,-4) -- (axis cs:-6.2832,4);
\draw[assint] (axis cs:-3.1416,-4) -- (axis cs:-3.1416,4);
\draw[assint] (axis cs:0.0000,-4) -- (axis cs:0.0000,4);
\draw[assint] (axis cs:3.1416,-4) -- (axis cs:3.1416,4);
\draw[assint] (axis cs:6.2832,-4) -- (axis cs:6.2832,4);
\end{groupplot}\end{tikzpicture}}{As seis funções trigonométricas. Seno, cosseno, secante e
cossecante têm período $2\pi$; tangente e cotangente, período $\pi$. Tracejado: assíntotas
verticais.}

Para $f(x)=A\sen\big(B(x-\alpha)\big)+C$ (ou com cosseno), $|A|$ é a \textbf{amplitude}\index{amplitude},
o período é $2\pi/|B|$, $\alpha$ é a translação horizontal (\emph{defasagem}) e $C$ é a
translação vertical.

\begin{exemplo}{1.8}{Esboçando uma senoide transformada}
Esboce $f(x)=3\sen\!\big(2(x-\tfrac\pi4)\big)+1$.
\solucao
Compressão horizontal por fator 2, defasagem de $\pi/4$ para a direita, alongamento
vertical por fator 3 e translação de 1 unidade para cima. Amplitude 3, período $\pi$.
\figtikz{\begin{tikzpicture}
\begin{axis}[painel,width=10cm,height=4.8cm,xmin=-4.9,xmax=7.2,ymin=-2.6,ymax=4.8,
  xtick={-4.712,-3.1416,-1.5708,0.7854,1.5708,3.1416,4.712,6.2832},xticklabels={$-\frac{3\pi}2$,$-\pi$,$-\frac\pi2$,$\frac\pi4$,$\frac\pi2$,$\pi$,$\frac{3\pi}2$,$2\pi$},ytick={-2,1,4}]
\draw[assint] (axis cs:-4.9,1) -- (axis cs:7.2,1);
\addplot[domain=-4.9:7.2,samples=300]{3*sin(deg(2*(x-pi/4)))+1};
\end{axis}\end{tikzpicture}}{Gráfico de $f(x)=3\sen(2(x-\pi/4))+1$: amplitude $3$, período $\pi$, deslocamento de fase $\frac\pi4$ e linha média $y=1$ (tracejada).}
\end{exemplo}

% ---------------------------------------------------------------------
\section{Funções inversas}

\subsection{Existência da inversa}

Uma função ``desfaz'' a outra se, ao aplicarmos uma e depois a outra, voltamos à entrada.

\begin{definicao}
Dada $f$ com domínio $D$ e imagem $R$, sua \textbf{função inversa}\index{funcao@função!inversa} (se existir) é a
função $f^{-1}$ com domínio $R$ e imagem $D$ tal que $f^{-1}(y)=x$ se $f(x)=y$. Ou seja,
\[
f^{-1}\big(f(x)\big)=x\ \text{ para todo } x\in D,\qquad
f\big(f^{-1}(y)\big)=y\ \text{ para todo } y\in R.
\]
\end{definicao}

\begin{nota}
$f^{-1}(x)$ não significa $1/f(x)$. A notação indica a função inversa.
\end{nota}

Para que $f^{-1}$ seja função, cada saída de $f$ deve vir de uma única entrada.

\begin{definicao}
$f$ é \textbf{injetora}\index{funcao@função!injetora} se $f(x_1)\neq f(x_2)$ sempre que $x_1\neq x_2$. Uma função tem
inversa se, e somente se, é injetora.
\end{definicao}

\begin{regra}{Regra: teste da reta horizontal}\index{teste da reta horizontal}
Uma função $f$ é injetora se, e somente se, toda reta horizontal corta o gráfico de $f$ no
máximo uma vez.
\end{regra}

\figtikz{\begin{tikzpicture}
\begin{groupplot}[group style={group size=2 by 1,horizontal sep=1.2cm},painel,width=5.6cm,height=4.8cm,xmin=-2.4,xmax=2.4,xtick={-1,1},ytick={2}]
\nextgroupplot[ymin=-0.5,ymax=4.4,title={(a) $x^2$}]
\draw[laranja] (axis cs:-2.4,2) -- (axis cs:2.4,2);
\addplot[domain=-2.1:2.1]{x^2};
\fill (axis cs:-1.414,2) circle (1.6pt); \fill (axis cs:1.414,2) circle (1.6pt);
\nextgroupplot[ymin=-4,ymax=4,title={(b) $x^3$}]
\draw[laranja] (axis cs:-2.4,2) -- (axis cs:2.4,2);
\addplot[domain=-1.6:1.6]{x^3};
\fill (axis cs:1.26,2) circle (1.6pt);
\end{groupplot}\end{tikzpicture}}{(a) $x^2$ não é injetora: falha no teste da reta
horizontal. (b) $x^3$ é injetora.}

\begin{estrategia}{encontrando a inversa}
\begin{enumerate}[label=\arabic*.]
\item Resolva $y=f(x)$ para $x$, em função de $y$.
\item Troque as variáveis e escreva $y=f^{-1}(x)$.
\end{enumerate}
\end{estrategia}

\begin{exemplo}{1.9}{Encontrando uma inversa}
Encontre a inversa de $f(x)=3x-4$.
\solucao
De $y=3x-4$ vem $x=\tfrac13y+\tfrac43$; logo $f^{-1}(x)=\tfrac13x+\tfrac43$, com domínio e
imagem $\R$. Verificação:
$f^{-1}(f(x))=\tfrac13(3x-4)+\tfrac43=x$.
\end{exemplo}

\textbf{Gráfico da inversa.} Se $(a,b)$ está no gráfico de $f$, então $(b,a)$ está no de
$f^{-1}$. Portanto o gráfico de $f^{-1}$ é a reflexão do gráfico de $f$ em relação à reta
$y=x$.

\figtikz{\begin{tikzpicture}
\begin{axis}[painel,width=6.2cm,height=6.2cm,xmin=-2.6,xmax=3,ymin=-2.6,ymax=3,xtick=\empty,ytick=\empty]
\addplot[cinza,dashed,domain=-2.6:3]{x} node[pos=0.97,below right,font=\scriptsize]{$y=x$};
\addplot[domain=-2:1.9,samples=100] ({x},{x^3/4+x/2}) node[pos=0.97,left,font=\scriptsize]{$f$};
\addplot[laranja,domain=-2:1.9,samples=100] ({x^3/4+x/2},{x}) node[pos=0.97,below,font=\scriptsize]{$f^{-1}$};
\fill[azul] (axis cs:1.2,1.032) circle (1.6pt) node[right,font=\scriptsize]{$(a,b)$};
\fill[laranja] (axis cs:1.032,1.2) circle (1.6pt) node[above left,font=\scriptsize]{$(b,a)$};
\draw[cinza,thin,dotted] (axis cs:1.2,1.032) -- (axis cs:1.032,1.2);
\end{axis}\end{tikzpicture}}{Se $(a,b)$ está no gráfico de $f$, então $(b,a)$ está no gráfico de $f^{-1}$: os gráficos são simétricos em relação à reta $y=x$.}

\textbf{Restrição de domínio.} Uma função que não é injetora pode ter inversa em um domínio
menor. Assim, $x^2$ restrita a $[0,\infty)$ tem inversa $\sqrt x$.

\begin{exemplo}{1.10}{Restringindo o domínio}
$f(x)=(x+1)^2$ não é injetora (é $x^2$ transladada 1 unidade para a esquerda). Em
$[-1,\infty)$ ela é injetora, com imagem $[0,\infty)$. Resolvendo $y=(x+1)^2$, vem
$x=-1\pm\sqrt y$; como $x\ge-1$, escolhemos o sinal $+$. Logo $f^{-1}(x)=-1+\sqrt x$, com
domínio $[0,\infty)$ e imagem $[-1,\infty)$.
\end{exemplo}

\subsection{Funções trigonométricas inversas}\index{funcoes trigonometricas inversas@funções trigonométricas inversas}

As funções trigonométricas não são injetoras, pois são periódicas. Restringindo o domínio,
obtemos as inversas.

\begin{definicao}
Para $-1\le x\le1$:
\[
\arcsen x=y \iff \sen y=x \text{ e } -\tfrac\pi2\le y\le\tfrac\pi2;\qquad
\arccos x=y \iff \cos y=x \text{ e } 0\le y\le\pi.
\]
Para todo $x$ real:
\[
\arctg x=y \iff \tg y=x \text{ e } -\tfrac\pi2<y<\tfrac\pi2;\qquad
\arccotg x=y \iff \cotg y=x \text{ e } 0<y<\pi.
\]
Para $|x|\ge1$, definem-se de modo análogo $\arccossec x$ (com
$-\tfrac\pi2\le y\le\tfrac\pi2$, $y\neq0$) e $\arcsec x$ (com $0\le y\le\pi$,
$y\neq\tfrac\pi2$). Em calculadoras e em livros em inglês aparecem as notações $\sin^{-1}$, $\cos^{-1}$, $\tan^{-1}$ (ou arcsin, arctan), com o mesmo significado.
\end{definicao}

\figtikz{\begin{tikzpicture}
\begin{groupplot}[group style={group size=3 by 1,horizontal sep=1.1cm},painel,width=5.2cm,height=5.2cm,
  xmin=-2,xmax=3.4,ymin=-2,ymax=3.4,
  xtick={-1,1},ytick={-1.5708,1.5708,3.1416},yticklabels={$-\frac\pi2$,$\frac\pi2$,$\pi$}]
\nextgroupplot[title={$y=\arcsen x$}]
\addplot[cinza,dotted,thin,domain=-2:3.4]{x};
\addplot[laranja!70,thin,domain=-1.5708:1.5708]{sin(deg(x))};
\addplot[domain=-1:1]{rad(asin(x))};
\nextgroupplot[title={$y=\arccos x$}]
\addplot[cinza,dotted,thin,domain=-2:3.4]{x};
\addplot[laranja!70,thin,domain=0:3.1416]{cos(deg(x))};
\addplot[domain=-1:1]{rad(acos(x))};
\nextgroupplot[title={$y=\arctg x$},xmin=-3.4,xtick={-2,2}]
\addplot[cinza,dotted,thin,domain=-3.4:3.4]{x};
\addplot[laranja!70,thin,domain=-1.25:1.25]{tan(deg(x))};
\addplot[domain=-3.4:3.4]{rad(atan(x))};
\draw[assint] (axis cs:-3.4,1.5708) -- (axis cs:3.4,1.5708);
\draw[assint] (axis cs:-3.4,-1.5708) -- (axis cs:3.4,-1.5708);
\end{groupplot}\end{tikzpicture}}{Cada inversa trigonométrica (em grafite) é a reflexão, em
relação à reta $y=x$ (pontilhada), da função trigonométrica restrita (em verde):
$\sen$ em $[-\frac\pi2,\frac\pi2]$, $\cos$ em $[0,\pi]$ e $\tg$ em
$(-\frac\pi2,\frac\pi2)$. As demais inversas são obtidas do mesmo modo.}

\begin{exemplo}{1.11}{Avaliando inversas trigonométricas}
(a) $\arcsen(-\sqrt3/2)$: o ângulo em $[-\pi/2,\pi/2]$ com seno $-\sqrt3/2$ é $-\pi/3$.

(b) $\arccos(\cos(5\pi/4))=\arccos(-\sqrt2/2)=3\pi/4$, e \emph{não} $5\pi/4$, pois
$\arccos$ toma valores em $[0,\pi]$.

(c) $\arcsen(\cos(2\pi/3))=\arcsen(-1/2)=-\pi/6$.
\end{exemplo}

% ---------------------------------------------------------------------
\section{Funções exponenciais e logarítmicas}

\subsection{Funções exponenciais}

Se uma população começa com $P_0$ indivíduos e cresce à taxa anual $r$, após $t$ anos ela
tem $P(t)=P_0(1+r)^t$ indivíduos. Esse é um exemplo de \textbf{função exponencial}\index{funcao@função!exponencial}: uma
função da forma $f(x)=b^x$, com $b>0$ e $b\neq1$. Seu domínio é $\R$ e sua imagem é
$(0,\infty)$. Se $b>1$, $b^x$ é crescente; se $0<b<1$, é decrescente. A exponencial
cresce muito mais rapidamente que qualquer potência: $2^x$ supera $x^2$ para $x$ grande.

\figtikz{\begin{tikzpicture}
\begin{axis}[painel,width=7cm,height=5cm,xmin=-3,xmax=3,ymin=-0.3,ymax=6.2,xtick={-1,1},ytick={1,2,4}]
\addplot[domain=-3:1.29]{4^x}; \addplot[domain=-3:2.58]{2^x};
\node[azul,font=\scriptsize,anchor=west] at (axis cs:1.35,5.5) {$4^x$};
\node[azul,font=\scriptsize,anchor=west] at (axis cs:2.2,4.2) {$2^x$};
\addplot[laranja,domain=-1.29:3]{(1/4)^x}; \addplot[laranja,domain=-2.58:3]{(1/2)^x};
\node[laranja,font=\scriptsize,anchor=east] at (axis cs:-1.35,5.5) {$(\frac14)^x$};
\node[laranja,font=\scriptsize,anchor=east] at (axis cs:-2.2,4.2) {$(\frac12)^x$};
\fill (axis cs:0,1) circle (1.5pt);
\end{axis}\end{tikzpicture}}{$b^x$ é crescente se $b>1$ (grafite) e decrescente se $0<b<1$ (verde). Todas passam por $(0,1)$.}

\begin{regra}{Regra: leis dos expoentes}\index{expoentes, leis dos}
Para $a>0$, $b>0$ e quaisquer $x$, $y$:
\[
b^xb^y=b^{x+y},\quad \frac{b^x}{b^y}=b^{x-y},\quad (b^x)^y=b^{xy},\quad
(ab)^x=a^xb^x,\quad \frac{a^x}{b^x}=\Big(\frac ab\Big)^x.
\]
\end{regra}

\subsection{O número $e$}

Se investimos $P$ reais a juros anuais $r$ capitalizados $n$ vezes por ano, após $t$ anos
temos $P(1+r/n)^{nt}$. Quando $n$ cresce (capitalização cada vez mais frequente), a
quantidade $(1+1/m)^m$ se aproxima de um número fixo,
\[
e\approx2{,}718282,
\]
e o montante tende a $Pe^{rt}$ (\emph{capitalização contínua}). O número $e$ é irracional.
A função $f(x)=e^x$ é a \textbf{exponencial natural}\index{exponencial natural}. Ela é a única exponencial cujo
gráfico tem reta tangente de inclinação exatamente 1 em $x=0$, fato que tornará as
derivadas e integrais de $e^x$ especialmente simples.

\subsection{Funções logarítmicas}

Como $b^x$ é injetora, ela tem inversa: o \textbf{logaritmo na base $b$},
\[
\log_b x=y \iff b^y=x,\qquad x>0 .
\]
Seu domínio é $(0,\infty)$ e sua imagem é $\R$. Escrevemos $\ln x=\log_e x$
(\textbf{logaritmo natural}\index{logaritmo natural}). Como inversas, $b^{\log_b x}=x$ para $x>0$ e
$\log_b(b^x)=x$ para todo $x$; em particular, $e^{\ln x}=x$ e $\ln(e^x)=x$.

\figtikz{\begin{tikzpicture}
\begin{axis}[painel,width=6.2cm,height=6.2cm,xmin=-3,xmax=4,ymin=-3,ymax=4,xtick={1},ytick={1}]
\addplot[cinza,dashed,domain=-3:4]{x} node[pos=0.97,below right,font=\scriptsize]{$y=x$};
\addplot[domain=-3:1.38]{exp(x)} node[pos=0.95,left,font=\scriptsize]{$e^x$};
\addplot[laranja,domain=0.05:4,samples=150]{ln(x)} node[pos=0.95,below,font=\scriptsize]{$\ln x$};
\end{axis}\end{tikzpicture}}{$y=e^x$ e $y=\ln x$ são inversas: seus gráficos
são simétricos em relação a $y=x$.}

\begin{regra}{Regra: propriedades dos logaritmos}\index{logaritmo!propriedades}
Se $a,b,c>0$, $b\neq1$ e $r\in\R$:
\[
\log_b(ac)=\log_b a+\log_b c,\qquad
\log_b\Big(\frac ac\Big)=\log_b a-\log_b c,\qquad
\log_b(a^r)=r\log_b a .
\]
\end{regra}

\begin{exemplo}{1.12}{Equações exponenciais e logarítmicas}
(a) $5^x=2$. Aplicando $\ln$: $x\ln5=\ln2$, logo $x=\ln2/\ln5$.

(b) $e^x+6e^{-x}=5$. Multiplicando por $e^x$: $(e^x)^2-5e^x+6=0$, isto é,
$(e^x-3)(e^x-2)=0$. Logo $x=\ln3$ ou $x=\ln2$.

(c) $\log_{10}\sqrt x+\log_{10}x=2$. O lado esquerdo é
$\log_{10}x^{3/2}=\tfrac32\log_{10}x$; logo $\log_{10}x=4/3$ e $x=10^{4/3}=10\sqrt[3]{10}$.
\end{exemplo}

\begin{regra}{Regra: mudança de base}\index{logaritmo!mudanca de base@mudança de base}
Para $a,b>0$ com $a,b\neq1$:
\[
a^x=b^{x\log_b a}\ \ (\text{em particular } a^x=e^{x\ln a}),\qquad
\log_a x=\frac{\log_b x}{\log_b a}\ \ (\text{em particular } \log_a x=\frac{\ln x}{\ln a}).
\]
\end{regra}

Graças a essas fórmulas, basta saber derivar e integrar $e^x$ e $\ln x$ para lidar com
qualquer base.

\begin{exemplo}{1.13}{A escala Richter}
Na escala Richter, de base 10, dois terremotos de magnitudes $R_1$ e $R_2$ e amplitudes
$A_1$ e $A_2$ satisfazem
\[
R_1-R_2=\log_{10}\Big(\frac{A_1}{A_2}\Big).
\]
Assim, um terremoto de magnitude 8 é 10 vezes mais intenso que um de magnitude 7, e
100 vezes mais intenso que um de magnitude 6. Comparemos o terremoto do Japão (2011,
magnitude 9) com o do Haiti (2010, magnitude 7,3).
\solucao
$9-7{,}3=\log_{10}(A_1/A_2)$, logo $A_1/A_2=10^{1{,}7}\approx50$. O terremoto do Japão
foi cerca de 50 vezes mais intenso.
\end{exemplo}

\subsection{Funções hiperbólicas}\index{funcoes hiperbolicas@funções hiperbólicas}

Certas combinações de $e^x$ e $e^{-x}$ aparecem com frequência. A forma de um cabo ou
corrente suspensa pelas extremidades, sob o próprio peso (a \emph{catenária}), é descrita
por uma delas.

\begin{definicao}
\[
\cosh x=\frac{e^x+e^{-x}}2,\qquad \senh x=\frac{e^x-e^{-x}}2,\qquad
\tgh x=\frac{\senh x}{\cosh x}=\frac{e^x-e^{-x}}{e^x+e^{-x}},
\]
\[
\cossech x=\frac1{\senh x},\qquad \sech x=\frac1{\cosh x},\qquad \cotgh x=\frac{\cosh x}{\senh x}.
\]
\end{definicao}

O nome vem da identidade $\cosh^2t-\senh^2t=1$: o ponto $(\cosh t,\senh t)$ percorre a
hipérbole $x^2-y^2=1$, assim como $(\cos t,\sen t)$ percorre o círculo $x^2+y^2=1$.

\figtikz{\begin{tikzpicture}
\begin{groupplot}[group style={group size=3 by 2,horizontal sep=1.1cm,vertical sep=1.3cm},painel,width=5.2cm,height=4.4cm,
  xmin=-3,xmax=3,xtick={-2,2}]
\nextgroupplot[title={$y=\cosh x$},ymin=-0.5,ymax=5,ytick={1,4}]
\addplot[domain=-3:3]{cosh(x)};
\addplot[cinza,dashed,thin,domain=-3:3]{exp(x)/2};
\addplot[cinza,dashed,thin,domain=-3:3]{exp(-x)/2};
\nextgroupplot[title={$y=\senh x$},ymin=-5,ymax=5,ytick={-4,4}]
\addplot[domain=-3:3]{sinh(x)};
\addplot[cinza,dashed,thin,domain=-3:3]{exp(x)/2};
\addplot[cinza,dashed,thin,domain=-3:3]{-exp(-x)/2};
\nextgroupplot[title={$y=\tgh x$},ymin=-2,ymax=2,ytick={-1,1}]
\addplot[domain=-3:3]{tanh(x)};
\draw[assint] (axis cs:-3,1) -- (axis cs:3,1);
\draw[assint] (axis cs:-3,-1) -- (axis cs:3,-1);
\nextgroupplot[title={$y=\cotgh x$},ymin=-4,ymax=4,ytick={-1,1}]
\addplot[domain=-3:-0.25]{1/tanh(x)};
\addplot[domain=0.25:3]{1/tanh(x)};
\draw[assint] (axis cs:-3,1) -- (axis cs:3,1);
\draw[assint] (axis cs:-3,-1) -- (axis cs:3,-1);
\nextgroupplot[title={$y=\operatorname{sech} x$},ymin=-0.5,ymax=1.5,ytick={1}]
\addplot[domain=-3:3]{1/cosh(x)};
\nextgroupplot[title={$y=\operatorname{csch} x$},ymin=-4,ymax=4,ytick={-2,2}]
\addplot[domain=-3:-0.25]{1/sinh(x)};
\addplot[domain=0.25:3]{1/sinh(x)};
\end{groupplot}\end{tikzpicture}}{As funções hiperbólicas. Em cinza tracejado,
$\pm\frac12e^{\pm x}$: para $|x|$ grande, $\cosh x$ e $\senh x$ se comportam como
$\frac12e^{|x|}$. Em verde, as assíntotas $y=\pm1$ de $\tgh$ e $\cotgh$.}

\begin{regra}{Regra: identidades hiperbólicas}\index{identidades hiperbolicas@identidades hiperbólicas}
\begin{tabular}{@{}ll@{}}
$\cosh(-x)=\cosh x$, \quad $\senh(-x)=-\senh x$ & $\cosh x+\senh x=e^x$, \quad $\cosh x-\senh x=e^{-x}$\\[2pt]
$\cosh^2x-\senh^2x=1$ & $1-\tgh^2x=\sech^2x$, \quad $\cotgh^2x-1=\cossech^2x$\\[2pt]
\multicolumn{2}{@{}l}{$\senh(x\pm y)=\senh x\cosh y\pm\cosh x\senh y$,\quad
$\cosh(x\pm y)=\cosh x\cosh y\pm\senh x\senh y$}
\end{tabular}
\end{regra}

\begin{exemplo}{1.14}{Avaliando funções hiperbólicas}
Se $\senh x=3/4$, então $\cosh^2x=1+(3/4)^2=25/16$; como $\cosh x\ge1$, temos
$\cosh x=5/4$. Daí $\tgh x=3/5$, $\cossech x=4/3$, $\sech x=4/5$ e $\cotgh x=5/3$.
\end{exemplo}

\textbf{Inversas hiperbólicas.} Como as hiperbólicas são combinações de exponenciais, suas
inversas se escrevem com logaritmos:
\[
\arcsenh x=\ln\!\big(x+\sqrt{x^2+1}\big),\qquad
\arccosh x=\ln\!\big(x+\sqrt{x^2-1}\big)\ (x\ge1),\]\[
\arctgh x=\frac12\ln\Big(\frac{1+x}{1-x}\Big)\ (|x|<1).
\]

% ---------------------------------------------------------------------
\section*{Resumo do capítulo}\markright{Resumo do capítulo}
\addcontentsline{toc}{section}{Resumo do capítulo}
\begin{itemize}
\item Uma função associa a cada entrada exatamente uma saída. Sem domínio explícito, o
  domínio é o conjunto dos $x$ em que a fórmula define um número real.
\item Teste da reta vertical: um gráfico de função é cortado no máximo uma vez por cada
  reta vertical. Uma função tem quantos zeros quiser, mas no máximo um intercepto $y$.
\item Para definir $g\circ f$, a imagem de $f$ deve estar contida no domínio de $g$; em
  geral $f\circ g\neq g\circ f$.
\item Funções pares são simétricas em relação ao eixo $y$; ímpares, em relação à origem.
\item A inclinação $m=\Delta y/\Delta x$ é a taxa de variação constante de uma função afim.
\item Translações, alongamentos, compressões e reflexões transformam gráficos conhecidos.
\item Ângulos em radianos: $\pi$ rad $=180^\circ$. Seno e cosseno têm período $2\pi$;
  tangente, período $\pi$.
\item Uma função tem inversa se, e somente se, é injetora; o gráfico de $f^{-1}$ é o reflexo
  do de $f$ em $y=x$. Restringindo o domínio, obtêm-se as inversas trigonométricas.
\item $b^x$ e $\log_b x$ são inversas; $\ln x=\log_e x$. Qualquer base se reduz à base $e$.
\item $\cosh x=(e^x+e^{-x})/2$, $\senh x=(e^x-e^{-x})/2$ e $\cosh^2x-\senh^2x=1$.
\end{itemize}

% Gerado por ferramentas/gerar.py; não edite à mão.
\section*{Exercícios selecionados}
\addcontentsline{toc}{section}{Exercícios selecionados}
Problemas escolhidos para desenvolver a resolução de problemas. As soluções completas,
com animações dos gráficos, estão no repositório do livro: \url{https://github.com/AleffAzevedo/calculo-1-pt/tree/main/exercicios/cap1}.

\begin{enumerate}[label=\textbf{\thechapter.\arabic*}, leftmargin=2.6em, itemsep=6pt]
\item \textbf{Composição e domínio.} Sejam $f(x)=\dfrac{3}{2x+1}$ e $g(x)=\dfrac{2}{x}$. Encontre $(f\circ g)(x)$ e
$(g\circ f)(x)$ e determine o domínio de cada composta.
\item \textbf{Custo e ponto de equilíbrio.} Uma confeitaria vende cupcakes em festivais. Montar a barraca custa R\$\,125,00 e cada
cupcake custa R\$\,0,75 para ser produzido.
(a) Escreva o custo total $C$ em função do número $x$ de cupcakes.
(b) Quanto custa produzir 160 cupcakes?
(c) Vendendo cada cupcake a R\$\,1,50, quantos ela precisa vender para começar a ter lucro?
\item \textbf{Identidade trigonométrica.} Verifique a identidade $\dfrac{1}{1-\sin\alpha}+\dfrac{1}{1+\sin\alpha}=2\sec^2\alpha$.
\item \textbf{Horas de luz do dia.} O número de horas de luz do dia em uma cidade é modelado por
$N(t)=12+3\sin\!\left[\frac{2\pi}{365}(t-79)\right]$, em que $t$ é o número de dias após
1º de janeiro.
(a) Encontre a amplitude e o período.
(b) Quantas horas de luz há no dia mais longo do ano? E no mais curto?
(c) Quantas horas de luz há 90 dias após 1º de janeiro?
\item \textbf{Custo de remoção de uma toxina.} O custo (em milhares de reais) para remover $p$ partes por bilhão (ppb) de uma toxina de um
lago é $C(p)=\dfrac{75p}{85-p}$, válido para $0\le p<85$.
(a) Calcule o custo para remover 25, 40 e 50 ppb.
(b) Encontre a função inversa $p=C^{-1}(C)$.
(c) Quanto da toxina se remove com R\$\,50 mil?
\item \textbf{Profundidade da maré.} A profundidade (em pés) da água em um cais é $D(t)=5\sin\!\left(\frac{\pi}{6}t-\frac{7\pi}{6}\right)+8$,
em que $t$ é o número de horas após a meia-noite. Determine o primeiro instante após a
meia-noite em que a profundidade é 11,75 pés.
\item \textbf{Expressões hiperbólicas.} Reescreva em termos de exponenciais e simplifique:
(a) $2\cosh(\ln x)$; (b) $\cosh 4x+\sinh 4x$; (c) $\cosh 2x-\sinh 2x$;
(d) $\ln(\cosh x+\sinh x)+\ln(\cosh x-\sinh x)$.
\item \textbf{Decaimento do iodo-131.} O iodo-131 decai segundo $Q(t)=Q_0e^{-0{,}08664t}$, com $t$ em dias. Quanto tempo leva
(arredondando para o dia mais próximo) para 95\% de uma amostra decair?
\item \textbf{Crescimento de bactérias.} Uma colônia de bactérias dobra a cada 12 horas e começa com 1000 bactérias.
(a) Usando $Q=Q_0e^{kt}$, determine $k$ (quatro casas decimais).
(b) Aproximadamente quanto tempo leva para haver 200\,000 bactérias?
\item \textbf{Escala Richter.} O terremoto de 1906 em São Francisco teve magnitude 8,3 na escala Richter. No mesmo
período, um terremoto de magnitude 4,9 no Japão causou poucos danos. Quantas vezes o de
São Francisco foi mais intenso?
\item \textbf{População sazonal.} A população de uma cidade litorânea, em milhares, é modelada por
$P(t)=82{,}5-67{,}5\cos\!\left(\frac\pi6t\right)$, com $t$ em meses ($t=0$ é 1º de janeiro).
Em que intervalos do ano a população é menor que 20 mil? E maior que 140 mil?
\end{enumerate}

\subsection*{Respostas}
{\small\begin{enumerate}[label=\textbf{\thechapter.\arabic*}, leftmargin=2.6em, itemsep=2pt]
\item $(f\circ g)(x)=\dfrac{3x}{x+4}$, com $x\neq0,-4$;\quad
$(g\circ f)(x)=\dfrac{2(2x+1)}{3}$, com $x\neq-\tfrac12$.
\item (a) $C(x)=0{,}75x+125$; (b) R\$\,245,00; (c) 167 cupcakes.
\item Some as frações: o denominador comum é $1-\sin^2\alpha=\cos^2\alpha$.
\item (a) Amplitude 3, período 365 dias; (b) 15 h e 9 h; (c) $\approx12{,}56$ h.
\item (a) $\approx$ R\$\,31,25 mil, R\$\,66,67 mil e R\$\,107,14 mil;
(b) $p=\dfrac{85C}{C+75}$; (c) 34 ppb.
\item $t=7+\frac6\pi\sin^{-1}(0{,}75)\approx8{,}62$ h, por volta de 8h37.
\item (a) $x+\frac1x$ ($x>0$); (b) $e^{4x}$; (c) $e^{-2x}$; (d) $0$.
\item $t=\dfrac{\ln 20}{0{,}08664}\approx35$ dias.
\item (a) $k=\frac{\ln2}{12}\approx0{,}0578$ por hora; (b) $\approx92$ horas.
\item $10^{3{,}4}\approx2512$ vezes.
\item $P<20$ para $t\in[0;\,0{,}74)\cup(11{,}26;\,12]$ (início de dez. a fim de jan.);
$P>140$ para $t\in(4{,}95;\,7{,}05)$ (fim de maio ao início de ago.).
\end{enumerate}}

